\begin{abstract}
小语言模型在多步知识查询中需要根据自身生成的历史继续求解。额外的合法操作既可能改变后续决策上下文，也会占用中间结果句柄，使直接接入参考程序后缀产生依赖错位。本文提出面向可执行知识问答的依赖一致恢复监督构造方法：由固定初始模型采集训练失败，保留共同成功前缀与一次实际分歧，通过句柄重映射接入参考后缀，并以执行重放核验训练目标。模型读取保留的偏离历史，损失集中于后续助手动作。基于Qwen2.5-3B-Instruct与LoRA，在相同初始化、成功材料、问题顺序及监督token预算下比较普通续训与恢复续训。在KQA Pro的2000题项目留出集上，两个配对续训种子的平均准确率由78.75\%提高至83.50\%，增加4.75个百分点，题级配对bootstrap的95\%区间为$[3.53,5.98]$。恢复组减少未完成回合和无效调用，相对配对分步基线的累计推理token分别下降8.84\%和4.17\%。保存训练历史上的另行离线构造诊断进一步揭示了原编号直接拼接的依赖失配。结果支持在确定性只读查询环境中利用参考程序构造恢复监督。
\end{abstract}

\noindent\textbf{关键词：}可执行知识问答；小语言模型；恢复监督；依赖重映射；执行反馈

\section{引言}

可执行知识问答将自然语言问题转换为知识库上的组合操作，例如实体定位、关系遍历、限定条件筛选以及计数或比较。中间结果同时承担下一步决策依据和后续操作输入的作用。KQA Pro提供问题、参考程序与答案，为研究操作组合提供了显式监督\cite{cao2022kqapro}。分步交互使语言模型能够读取执行观察\cite{yao2023react}，也要求它持续管理已经产生的实际历史。

成功演示通常呈现沿参考路径展开的上下文。模型自由运行时，可能增加一次无关查询、得到空集合或反复选择无效依赖，即使正确后续操作仍可执行，也可能在已有历史中停止推进。采用追加式句柄管理中间结果时，一次额外成功操作还会改变随后结果的物理编号。若直接复制参考后缀，后续依赖可能指向错误对象。

本文研究在已有参考程序与确定性查询执行器的条件下，怎样从模型实际访问的偏离历史构造后续监督。方法保留最终失败回合中的首次参考分歧，通过句柄映射连接正确后缀，执行验证后将其作为训练目标。实际分歧及反馈作为上下文，后续动作作为监督内容，核心是保持参考依赖结构与当前执行状态之间的一致对应。

为评价该训练构造，本文建立共享初始检查点、普通成功材料和逻辑后缀目标的配对对照。两组区别在于后缀之前是否保留真实分歧及反馈，各自依赖编号适应其实际状态，监督token数量逐记录匹配。固定项目留出集上的平均增量为4.75个百分点；执行状态与成本分析描述性能变化的表现，离线构造诊断则检验句柄错位对参考目标的影响。

本文贡献包括：构建结合真实偏离历史、依赖映射和重放核验的恢复目标生成方法；建立保留完整所选历史且匹配目标token的配对监督方案；以问答准确率、完成状态、推理成本和后缀构造有效性形成相互衔接的实证分析。研究环境为知识库固定、操作只读、结果可保留且参考后缀可重放的程序式查询。

\section{相关工作}

\paragraph{策略访问状态与反馈利用。}
DAgger通过聚合策略实际访问状态上的专家动作，处理专家演示与实际执行状态分布之间的差异\cite{ross2011dagger}。ReAct将推理、动作和观察组织为交互过程\cite{yao2023react}；Self-Refine利用自反馈迭代修改输出\cite{madaan2023selfrefine}，Reflexion使用语言反馈和记忆支持后续尝试\cite{shinn2023reflexion}。这些工作为状态覆盖和反馈组织提供基础，本文进一步关注给定程序依赖结构的训练目标构造。

\paragraph{恢复轨迹训练。}
Agent-R以蒙特卡洛树搜索形成成功与失败分支，并通过修订轨迹开展迭代自训练\cite{yuan2025agentr}。STeP由教师提供反思与纠正，以部分掩码控制目标位置\cite{chen2025step}。OEC在学生历史中切换到专家策略，训练经环境筛选的专家续写\cite{lauffer2025oec}。ReTerm利用更强教师从失败终态及历史继续完成任务，并进一步开展在线恢复学习\cite{xie2026reterm}。这些方法共享学生上下文与后续改进动作的组合；本文采用已有参考程序提供操作，通过显式句柄映射将其连接到保留的历史。

\paragraph{纠正行为的学习目标。}
SCoRe以多轮强化学习训练自我纠正\cite{kumar2025score}，Fission-GRPO把执行错误转换为恢复任务并在强化学习中重新采样\cite{zhang2026fission}。本文的主要对照采用监督续训，差异集中于纠正目标来源、执行结构和监督条件。表\ref{tab:related-zh}比较方法设计；性能实验采用统一任务内的配对控制。

\begin{table}[htbp]
\centering\small
\caption{恢复监督的目标来源与构造条件}
\label{tab:related-zh}
\RecoveryRelatedTableZH
\end{table}

\section{依赖一致的恢复监督方法}

\subsection{执行环境与保留历史}

设固定知识库为$G$，问题为$q$，执行状态为$s_t$，模型可见历史为$H_t$。查询动作$a_t=(f_t,\ell_t,D_t)$包含函数、字面参数和依赖句柄列表，执行过程为
\begin{equation}
 (s_{t+1},o_t)=\mathcal E(G,s_t,a_t),\qquad
 H_{t+1}=H_t\mathbin{\Vert}(a_t,o_t).
\end{equation}
成功操作追加结果句柄，被拒绝的动作返回反馈而不新增成功结果；结束动作提交指定句柄。空结果可由合法查询产生，回合正确性由完整答案评价。

训练问题具有参考程序$P^*$和答案$y^*$。构造器执行参考程序获得动作与观察，同时回放初始模型$I$的实际轨迹，核对消息、事件和提交结果。对最终未正确作答的回合，寻找工具名称、参数和执行成功状态均与参考一致的最长共同前缀，其长度为$k\geq1$；随后实际动作及观察构成分歧$d=(\hat a_{k+1},\hat o_{k+1})$。

参考分歧用于确定构造起点。由于同一问题可能存在多条有效程序，样本准入同时结合完整回合失败以及正确后缀的实际可执行性。其目标是在保留真实历史后继续完成原问题。

\subsection{依赖映射与执行核验}

设$m(i)$表示参考结果$i$在当前回合中的对应句柄。重放共同前缀时建立映射，对后续参考动作$a_i^*=(f_i^*,\ell_i^*,D_i^*)$，构造
\begin{equation}
 \rho_m(a_i^*)=
 \left(f_i^*,\ell_i^*,[m(j):j\in D_i^*]\right).
 \label{eq:remap-zh}
\end{equation}
动作执行得到句柄$u$后，更新$m(i)=u$；结束操作采用相同映射。分歧产生的结果保留在回合中，参考对应关系由共同成功前缀和重新执行的后缀结果建立。

依赖不变量为：后续引用均来自共同前缀的已验证结果，或已经执行的参考后缀步骤。固定知识库、确定性只读操作及保留的中间结果，使该对应能够在续行中维持。构造器检查依赖可用性和接口类型，执行映射后的动作，将返回句柄规范到对应编号后比较参考观察，并核验最终答案；全部操作满足原调用预算。

例如，共同前缀产生$h_0$，参考后缀原本依次产生$h_1$、$h_2$并提交$h_2$。若额外查询已占用$h_1$，第一项后缀操作将产生$h_2$，下一项应引用$h_2$、产生$h_3$并提交$h_3$。保留原编号则可能引用额外查询或前一项后缀结果。映射调整这些依赖，函数与字面参数保持参考程序的定义。

\begin{figure}[htbp]
\begin{center}\small
\fbox{\parbox{0.92\linewidth}{\centering 训练问题、参考程序与固定初始模型 I}}
\par\vspace{0.3em}$\downarrow$\par\vspace{0.3em}
\fbox{\parbox{0.92\linewidth}{\centering 真实失败轨迹：共同成功前缀与首次参考分歧}}
\par\vspace{0.5em}
\begin{minipage}[t]{0.46\linewidth}\centering
\fbox{\parbox{0.90\linewidth}{\centering\textbf{普通续训 A}\\
共同前缀\\$\downarrow$\\执行参考后缀}}
\end{minipage}\hfill
\begin{minipage}[t]{0.46\linewidth}\centering
\fbox{\parbox{0.90\linewidth}{\centering\textbf{恢复续训 C}\\
共同前缀＋真实分歧及反馈\\$\downarrow$\\依赖映射后执行参考后缀}}
\end{minipage}
\par\vspace{0.4em}$\downarrow$\par\vspace{0.3em}
\fbox{\parbox{0.92\linewidth}{\centering 核验观察与答案；仅监督后缀助手动作\\两侧共享成功材料、问题顺序与监督 token 预算}}
\end{center}
\caption{离线恢复监督构造。参考程序用于训练目标生成，正式评价由模型自由执行。}\label{fig:method-zh}
\end{figure}


\subsection{配对语料与监督目标}

令$S^*$为包括提交操作的参考后缀。普通续训组A直接从共同前缀继续，恢复组C先保留实际分歧，两类记录为
\begin{align}
 x^A&=(q,H_k,S^*,\mathrm{Done}),\\
 x^C&=(q,H_k,d,\rho_m(S^*),\mathrm{Done}).
\end{align}
两组均在动作之间插入实际工具观察；后缀具有相同的逻辑函数、字面参数与依赖图，物理句柄数字可不同。

对序列$x_{1:L}$定义监督掩码$w_t$，采用因果语言建模损失
\begin{equation}
 \mathcal L(\theta)=-\frac{1}{\sum_x\sum_t w_t}
 \sum_x\sum_t w_t\log p_\theta(x_t\mid x_{<t}).
\end{equation}
配对记录中的问题、共同前缀、分歧及全部工具观察不参与损失，后续助手动作和结束文本作为监督。两组另混合完全相同的普通成功轨迹，以保持常规查询能力。

每组续训总预算为788534个监督token，普通成功材料和配对后缀各占394267。每一配对取两侧可用目标数与剩余预算的最小值，对超额目标从尾部屏蔽，完整上下文保留。每组最终包含5772条记录、2694个不同问题；配对部分由418个不同配对重复采样形成，两侧问题次序、重复次数及逐记录目标数一致。

构造流程为：
\begin{enumerate}
\item 回放初始模型的训练回合并执行参考程序，选择最终失败、记录一致且参考答案有效的样本。
\item 定位非空共同成功前缀及其后的实际分歧。
\item 为两组分别恢复前缀，恢复组追加真实分歧，建立参考结果到当前句柄的映射。
\item 在调用预算内重映射并执行参考后缀，核对观察及最终答案。
\item 保存配对，设置后缀监督掩码，在语料构建时匹配目标token。
\end{enumerate}
全部排除原因随记录保存。参考后缀用于离线构造，正式评价由模型依据问题和实际执行观察自行生成动作。

\section{实验设置}

\subsection{数据划分与对照}

实验使用KQA Pro及其规则执行器\cite{cao2022kqapro}，从官方训练数据形成5000题训练、500题开发和2000题项目留出划分。训练和开发沿用较早固定选择；新增留出部分排除8500个历史问题／程序分组，并另行排除与官方验证集精确匹配的分组。三部分之间的精确问题和程序交集为零。

初始模型生成5000条训练轨迹，4998条回放一致，其中4339条正确、659条失败。最终准入418个配对，占回放有效失败的63.43\%；另240条缺少共同成功前缀，1条参考答案不匹配。418次保留分歧中，414次可执行、4次被拒绝，主要覆盖合法偏离之后的继续求解。

主要对照为A与C：共享初始模型I和匹配的续训语料，续训种子分别为20261003、20261004。一次性程序生成模型P作为求解形式及成本参照，在4999条程序示例上训练三遍，每遍273426个监督token。P的监督形式与预算不同，方法增量以配对A/C比较为依据。

\subsection{模型与推理参数}

基座采用Qwen2.5-3B-Instruct\cite{qwen2024qwen25}，LoRA秩16、缩放32、dropout 0.05，适配实现指定的全部线性层\cite{hu2022lora}。5000个训练问题中4999个参考执行得到目标答案，I使用其成功分步演示训练两遍，每遍788534个监督token。A/C各续训一遍，AdamW学习率$10^{-4}$，每设备批大小2、梯度累积8，采用BF16、梯度检查点与SDPA。

分步推理采用贪心解码、批大小8、8192上下文上限；每题最多24次工具调用、2048个生成token，每轮最多192个生成token。全部问题保留在评价分母。五个检查点、推理配置、反馈形式、主估计与成本定义均在留出答案评分前固定，已有执行审计核对10000条保存结果，回放差异为零。

\subsection{指标与配对统计}

主指标为完整答案准确率。设$y_{i,s}^{M}\in\{0,1\}$为问题$i$、模型$M$、续训种子$s$的正确性，预登记主估计为
\begin{equation}
 \widehat\Delta=\frac{100}{N}\sum_{i=1}^{N}
 \frac{(y_{i,1}^{C}-y_{i,1}^{A})+(y_{i,2}^{C}-y_{i,2}^{A})}{2},
 \qquad N=2000.
\end{equation}
进行5000次题级配对百分位bootstrap，同一问题的两个种子结果共同进入抽样，区间对应固定检查点下的问题不确定性。完成状态、无效调用及未完成／完成错误／正确之间的转移描述执行表现。推理成本累加每轮未填充的完整输入与生成token，包含重复历史输入和生成EOS。

\section{结果与分析}

\subsection{项目留出集问答效果}

\begin{table}[htbp]
\centering\small
\caption{同一2000题项目留出集的问答效果；增量与区间单位为百分点}
\label{tab:main-zh}
\RecoveryMainTableZH
\end{table}

表\ref{tab:main-zh}中，第一配对准确率由78.25\%提高至82.90\%，第二配对由79.25\%提高至84.10\%。两次增量分别为4.65与4.85个百分点，95\%区间为$[3.10,6.25]$和$[3.35,6.40]$；平均增量4.75个百分点，区间$[3.53,5.98]$。在固定推理预算下，恢复上下文及后缀监督的组合提高了完整问答准确率。

第一配对改正177题、退步84题，净增93题；第二配对改正171题、退步74题，净增97题。两种子点估计方向一致，同时保留改正与退步的完整计数。P达到76.15\%，两组C分别高出6.75与7.95个百分点，其交互成本差异在下节列出。

\subsection{完成状态与推理成本}

\begin{table}[htbp]
\centering\small
\caption{项目留出集上的分步执行与累计token成本}
\label{tab:cost-zh}
\RecoveryCostTableZH
\end{table}

未完成题由198降至80、由205降至91，均对应达到24次调用上限。无效调用由3725降至1268、由3845降至1216，分别减少65.96\%与68.37\%。状态转移中，未完成与正确之间的净转移贡献72与75题，双方均完成时净增21与22题；平均95道新增正确题中，73.5道跨越完成边界。部分新增完成仍然回答错误，已有正确题也存在退步，整体收益同时包含终止改善及完成回合内部的答案变化。

相对A，C累计输入与生成token合计减少8.84\%和4.17\%；其生成token分别增加3.88\%和16.35\%，总量下降主要来自重复输入减少。P的累计总token为804096，C的两次运行分别为其23.64与24.72倍。恢复监督在匹配分步基线上获得质量与累计token改进，一次性程序路线在该成本指标下仍更低。

\subsection{参考后缀构造诊断}

\begin{table}[htbp]
\centering\small
\caption{414个可严格比较历史训练配对的离线后缀构造；原418总记录中另4个单列为回放不可比}
\label{tab:construction-zh}
\RecoveryConstructionTableZH
\end{table}

另以全部418个历史准入配对开展CPU构造诊断。各处理恢复相同的已保存共同前缀及实际分歧：直接拼接保留原参考编号，映射执行调整依赖，完整构造再按参考一致性验收。直接拼接在调用被拒绝后继续尝试剩余固定列表，直至提交或达到原预算，期间不重新规划。

414个配对在本机可精确比较；另4个拒绝动作的\texttt{UnboundLocalError}详情在旧环境与Python 3.12中措辞不同，动作、异常类别与拒绝状态一致，按严格回放规则保留为不可比。可比较的414次分歧均新增了句柄。

直接拼接有65条后缀全部调用可执行，121条提交答案，9条答案与参考相同；映射执行的414条均完成后缀并得到一致答案。直接拼接的348条存在中间依赖错位，391条存在提交句柄错位，两类可重叠。9条答案相同的轨迹仍均有参考依赖来源差异及至少一次无效调用，其中5条还出现内部结果类型或内容差异。答案匹配因此需要与依赖来源共同用于训练目标验收。

映射未筛选与完整验收在414条上产生相同序列，验收没有新增剔除。这与样本已经历史准入的条件相符：关闭一个检查未必产生新的训练语料。该诊断刻画构造失效方式及目标一致性，学习效果则由前述留出A/C实验单独评价。

\subsection{反馈文字诊断}

辅助推理实验保留实际执行状态及全部成功观察，将失败时的具体错误类别和详情替换为通用文本。四个固定A/C检查点在500题开发集上的变化为$+0.2$至$+0.4$个百分点，区间均含零。该窄干预未显示详细错误文字的稳定增益；记录中另有5条轨迹在首次干预前已出现生成差异，原因尚未确定。训练分歧主要是可执行动作，后续行为分析应同时关注更完整的交互历史。

\section{讨论与适用范围}

本文形成两类相互补充的证据：受控续训实验从原问题开始评价学习后的策略，显示固定留出集上的正向增量；离线构造诊断则解释保留历史改变后，参考依赖怎样连接到实际结果。两者共同支持经验证的后缀监督构造，分别测量学习效果与构造有效性。

当前研究采用单一3B基座、一个知识查询基准及两个共享初始模型与语料的续训种子。项目留出控制精确问题／程序重复，知识库实体、关系和部分结构仍共享；题级区间条件于固定检查点。准入样本具有共同成功前缀和足够的重放预算，主要是合法偏离。固定只读查询满足现有状态假设，修改外部环境的工具还需要其他状态恢复机制。

相同正常与偏离历史下的模型续行评价已登记但尚未运行，它将进一步控制模型接手状态，补充当前自由运行及构造层证据。现有准确率增量对应完整恢复监督方案，各组成对学习效果的独立贡献需要具有明确处理差异的实验继续分析。

\section{结论}

本文提出面向可执行知识问答的依赖一致恢复监督构造，以真实失败历史提供上下文、参考程序提供后续目标，通过句柄重映射和执行重放连接二者。在匹配监督量的条件下，2000题KQA Pro项目留出集的平均准确率由78.75\%提高至83.50\%，未完成回合与相对分步基线的累计推理token减少。构造诊断进一步显示，陈旧句柄会破坏参考依赖，即使最终答值偶然相同也可能存在路径失配。该研究为小模型在结构化、可重放知识查询中的恢复学习提供了具体方法和实证结果。
